\documentclass[../../main.tex]{subfiles}
\begin{document}

\begin{center}
  {\fontsize{48pt}{58pt}\selectfont \textbf{2008}\normalsize\textbf{年度}}
\end{center}

\begin{table}[htb]
  \centering
  \caption{出題テーマと難易度}
  \begin{tabular}{|c|p{9cm}|c|}
    \hline
    \textbf{問題番号} & \textbf{テーマ} & \textbf{難易度} \\
    \hline
    第1問 & 対数曲線と直線が交わらない条件 & \\
    \hline
    第2問 & 正四面体上の移動確率 & \\
    \hline
    第3問 & 【乙】四直線と平行四辺形／【甲】二等辺三角形の角 & \\
    \hline
    第4問 & 絶対値を含む二曲線の共有点数 & \\
    \hline
    第5問 & 円柱を斜めの平面で切った体積 & \\
    \hline
    第6問 & 同緯度の経路と大円航路の長さ & \\
    \hline
  \end{tabular}
\end{table}

\begin{table}[htb]
  \centering
  \caption{解答}
  \renewcommand{\arraystretch}{1.4}
  \begin{tabular}{|c|c|p{9cm}|}
    \hline
    \textbf{問題番号} & \textbf{小問} & \textbf{答え} \\
    \hline
    第1問 & -- & $p>0$かつ$1+q+\log p>0$ \\
    \hline
    第2問 & -- & $1-3\left(\dfrac{2}{3}\right)^n+3\left(\dfrac{1}{3}\right)^n$ \\
    \hline
    第3問 & 【乙】 & 条件を満たす平面が存在（証明問題） \\
    \cline{2-3}
    & 【甲】 &  \\
    \hline
    第4問 & -- & $|a|<2$で2個，$|a|=2$または$|a|=3\sqrt{2}/2$で3個，その他の$|a|>2$で4個 \\
    \hline
    第5問 & -- & $3\sqrt{3}-\dfrac{4\pi}{3}$ \\
    \hline
    第6問 & -- & $R_2<0.97R_1$（証明問題） \\
    \hline
  \end{tabular}
\end{table}

\paragraph{試験形態}
6問．試験時間は要確認．

\end{document}
